Для формализации персонализации применяется связка из вероятностной оценки знания и алгоритма последовательного выбора действия. Вероятностная модель отслеживает состояние освоения навыков, а управляющий модуль максимизирует ожидаемую учебную пользу.

В проекте используется совместимая с прикладным контуром конфигурация: трекинг знаний на уровне навыков и стратегия выбора заданий с учетом ожидаемого прироста. Формально оптимизация может быть записана как:

\begin{equation}
\pi^* = \arg\max_{\pi} \mathbb{E}\left[\sum_{t=0}^{T} \gamma^t r_t\right],
\end{equation}

где $r_t$ отражает учебный прогресс обучающегося на шаге $t$.
